% !PRINTED 464
\weberpage{464}
%<<CONT
bewiesen haben, ist der Grad $n_1$ des Körpers $\Omega(\beta)$ ein
Theiler von $n$. Ist nun durch $\Omega(\beta)$ der Körper $\Omega'$
nicht erschöpft, so nehmen wir eine Grösse $\gamma$ in $\Omega'$,
aber nicht in $\Omega(\beta)$; dann ist der Körper $\Omega(\beta,
\gamma)$ ein Theiler von $\Omega'$ und von $\Omega(\alpha)$ und hat
seinerseits $\Omega(\beta)$ zum Theiler. Der Grad von $\Omega(\beta,
\gamma)$ ist also grösser als der von $\Omega(\beta)$ und kleiner
als der von $\Omega(\alpha)$. Ist damit $\Omega'$ noch nicht
erschöpft, so fahren wir so fort, müssen aber endlich zum Abschluss
kommen, da die Grade der Körper $\Omega(\beta), \Omega(\beta,
\gamma) \ldots$ immer wachsen und doch kleiner als $n$ bleiben.
Daraus folgt:

\gesperrt{5. Jeder Theiler von $\Omega(\alpha)$, der den Körper
$\Omega$ enthält, ist ein algebraischer Körper $\Omega(\beta)$ über
$\Omega$. Der Grad von $\Omega(\beta)$ ist ein Theiler des Grades
von $\Omega(\alpha)$, und wenn beide Körper den gleichen Grad haben,
so sind sie identisch.}

Wir können daher unsere Definition auch so fassen:

\gesperrt{6. Ein algebraischer Körper über $\Omega$ ist primitiv,
wenn er ausser $\Omega$ und sich selbst keinen Körper über $\Omega$
zum Theiler hat.}

\begin{center}
\S.\,145.

\gesperrt{Normalkörper. Galois'sche Resolvente.}
\end{center}

Ist $\Omega(\alpha)$ ein algebraischer Körper, so sind die
conjugirten Körper
% [sic] The source prints alpha_{m-1}; §144 defines the conjugate fields only through alpha_{n-1}, where n is the irreducible-factor degree. Preserved diplomatically.
\[
\Omega(\alpha), \Omega(\alpha_1) \ldots \Omega(\alpha_{m-1})
\]
alle von gleichem Grade, und wenn also einer von ihnen im anderen
enthalten ist, so sind beide identisch.

Ein Körper, der mit allen seinen conjugirten Körpern identisch
ist, heisst ein \gesperrt{Normalkörper}. In den Normalkörpern
herrschen viel einfachere Gesetze, und der grosse Fortschritt, den
die Algebra \gesperrt{Galois} verdankt, beruht im Wesentlichen
darauf, dass beliebige Körper auf Normalkörper zurückgeführt
werden. Die Normalkörper heissen daher auch \gesperrt{Galois'sche
Körper}.

Wenn ein Körper $\mu^{\mathrm{ten}}$ Grades $\Omega(\varrho)$ die
Eigenschaft hat, dass die zu $\varrho$ conjugirten Zahlen $\varrho_1,
\varrho_2 \ldots \varrho_{\mu-1}$ alle in $\Omega(\varrho)$
enthalten sind, so ist es ein Normalkörper; denn dann sind die
Körper $\Omega(\varrho_1), \Omega(\varrho_2) \ldots
\Omega(\varrho_{\mu-1})$ auch alle in $\Omega(\varrho)$ enthalten
und also alle mit $\Omega(\varrho)$ identisch.
